\section{Lineage and Context}
\label{sec:checkpoints-lineage}

\subsection{2016 Language-Model Lesson}
\label{sec:checkpoints-lineage-2016}

The research lineage begins with a 2016 language-model experiment whose finding was that local text
imitation does not conserve consequence. A model that reproduces the surface patterns of a corpus
does not reproduce the causal structure that produced the corpus. Predictions can be locally
plausible and globally incoherent. This lesson established the foundational distinction between
description and consequence that runs through the entire thesis: describing a process is not the
same as demonstrating that the process is sound.

The checkpoint ledger is a direct institutional response to this lesson. The ALIVE verdict mechanism
was designed specifically to prevent the research program from falling into the same trap at the level
of knowledge claims. A research report that \emph{describes} a doctrine does not constitute evidence
that the doctrine has been manufactured at minimum viable density. The gate criteria quantify the
minimum density. The audit verification script in RESEARCH\_CRITERIA.md can be re-run at any time to
confirm that the counts remain valid. The ledger converts the 2016 lesson from a philosophical caution
into an operational constraint on the research program itself.

\subsection{Enterprise Process Gap}
\label{sec:checkpoints-lineage-enterprise-gap}

The enterprise context that motivates process intelligence research is the gap between declared
processes and enacted processes in production environments. Organizations describe their workflows in
BPMN diagrams, SLA documents, and governance frameworks. Actual execution departs from these
declarations in ways that are invisible unless process-mining tools are applied to the execution
logs.

The checkpoint ledger addresses an analogous gap at the research program level. The declared
research process is the gate criteria table: twelve categories, each with a minimum count. The
enacted research process is the actual artifact inventory at the moment of assessment. The ledger
makes the gap visible by recording both the declared threshold and the actual count for each gate.
Where the count exceeds the threshold, the gate passes. Where it falls short, the gate fails and the
verdict is PARTIAL. The enterprise process gap is not merely a domain problem to be studied; it is
a pattern that the research program must not reproduce in its own operation.

\subsection{Relationship to the Chatman Equation}
\label{sec:checkpoints-lineage-chatman-equation}

The Chatman Equation $A = \mu(O^{*})$ expresses the thesis that any admissible process artifact is a
function of the organizational substrate $O^{*}$ applied through the manufacturing operator $\mu$.
In the checkpoint context, the admissibility boundary is stated formally as $R \vdash P_i = \mu(O^*, T, L)$,
where $R$ is the research program's type-law authority, $P_i$ is the process artifact being admitted,
$T$ is the type-law surface, and $L$ is the event log evidence. This invariant is listed as the first
mathematical invariant in PROCESS\_INTELLIGENCE\_ALIVE\_001.md and is marked verified at the time of
sealing.

The checkpoint system operationalizes the Chatman Equation by requiring that every ALIVE verdict be
traceable to a manufacturing event with a named witness. The admissibility boundary is not satisfied
by assertion; it is satisfied by the presence of artifacts whose counts can be verified against the
declared thresholds. The gate criteria table is the $O^{*}$ substrate encoded as a minimum density
specification. The artifact inventory is the $\mu$ application result. The ALIVE verdict is the
admission decision.

\subsection{Relationship to Receipts and Replay}
\label{sec:checkpoints-lineage-receipts-replay}

Receipts and replay are the two mechanisms by which process evidence acquires its authority in the
thesis architecture. A receipt is a typed, witnessed, bound evidence artifact: it records what was
manufactured, under which witness marker, and against which criterion. Replay is the ability to
re-execute a manufacturing pipeline and verify that the output matches the receipt. Together they form
the evidence chain that the checkpoint ledger references.

Every ALIVE verdict in the ledger cites at least one receipt. The RECEIPT\_REGISTRY.md indexes seven
canonical research program receipts: PAPER\_CANON\_RECEIPT, PM4PY\_ORACLE\_RECEIPT,
WASM4PM\_GAP\_RECEIPT, LIFECYCLE\_RECEIPT, MA\_RECEIPT, STANDARDS\_RECEIPT, and
ADVERSARIAL\_RECEIPT. Each receipt names the workflow that manufactured it, the witness law it
satisfies, the result type, the criteria met, and the gaps remaining. The SUBSTRATE\_COMPLETE\_001
checkpoint extends this pattern to BLAKE3-sealed generation receipts for five wasm4pm core modules:
mining (839 lines), conformance (647 lines), replay (467 lines), lifecycle (839 lines), and the Blue
River Dam orchestrator (629 lines). The receipt chain is the audit trail that makes ALIVE verdicts
replayable rather than merely asserted.
